\begin{document}

\section{実験設定}

\subsection{データとchain}

対象はATLASv2由来ログから構成した23 behavior chainである．各chainは，起点となる観測行，対象ホスト，focus process，time window，gold behavior steps，critical evidence slotsを持つ．旧27-chain実験から，現在の評価対象として妥当な23 chainを採用し，goldはCBC alert summaryだけを根拠にしない形へ統一した．これにより，Stage 3でalert summary rowsを隠した場合でも，残存telemetryから回答可能な範囲を評価できる．

chainは大きく3種類に分ける．明示的実行チェーンは，プロセス名，コマンドライン，ネットワーク接続，レジストリ操作が直接観測される行動であり，詳細分類では\code{network\_service\_behavior}と\code{command\_shell\_execution}を含む．複数段ツールチェーンは，tshark，batch，Sublime，cmd，Pythonなど複数のプロセスやツールをまたぐ行動であり，\code{collection\_or\_tool\_invocation}と\code{script\_execution\_chain}を含む．意味解釈型チェーンは\code{persistence\_registry\_run\_key}，すなわちDiscord Run Keyの1 chainであり，レジストリ永続化の意味づけを含むため，カテゴリ全体の結論ではなくcase studyとして扱う．

\subsection{3段階の入力条件}

表\ref{tab:stages}に，3つのstageを示す．Stage 1はCBC alert rowsを入力に含めるアラート支援条件である．Stage 2はhost，focus process，time windowのみを入力とし，DB内にはalert summary rowsを残す．Stage 3はStage 2と同じ初期入力だが，検索対象からCBC alert summary rowsを除外する．ただし，CBC EDR/NGAV event telemetry，Security，Sysmon，DNS，browser historyは残す．

\begin{table*}[t]
\centering
\scriptsize
\caption{3段階の入力条件}
\label{tab:stages}
\begin{tabular}{p{0.13\linewidth}p{0.26\linewidth}p{0.52\linewidth}}
\toprule
Stage & 入力 & 意味 \\
\midrule
Stage 1 & CBC alert + process/time & アラート支援の復元 \\
Stage 2 & process/timeのみ & DB内の要約行は利用可能 \\
Stage 3 & process/timeのみ & alert summary rowsをSQL検索対象から除外 \\
\bottomrule
\end{tabular}
\end{table*}

この設計により，初期入力の強さと，DB内の高密度なアラート要約行の有無を分けて観察できる．Stage 3はCBC全体を除外する条件ではない．CBC alert summary rowsのみを検索対象から外し，CBC EDR/NGAV telemetryは残す．したがって，本稿の結論は「alert summaryなしでも一定の行動復元が可能」という範囲に留め，「任意の低レベルログだけで十分」とは主張しない．

\subsection{実行と採点}

比較対象は\code{gpt-4.1-mini}，\code{gpt-5.4-mini}，および\code{gpt-5.5 low raw}である．各モデルについて23 chain，3 stage，3 set相当を評価し，各モデル207 rowsを得た．\code{gpt-4.1-mini}と\code{gpt-5.4-mini}の第三セットは，同一条件で新規実行したformal23 replicateではなく，旧27-chain実験から現在の23-chain対象だけを抽出した実質第三セットである．\code{gpt-5.5 low raw}は3セット完了しているが，要求したJSON/\code{code\_steps}出力契約に従わなかったため，raw-output salvageとして採点した．

採点はCodex内の複数レビューにより，component rubricで行った．各runは2つの独立レビューを基本とし，判定が割れた場合は第三レビューで裁定した．主指標は，action component recall，critical evidence recall，behavior sequence order，candidate claim precision，overclaim slot countである．CBC alert-summary-only evidenceはnon-alert critical evidenceとして数えない．\code{gpt-5.5 low raw}は内容包含に基づく実質的な復元能力を見るために採点対象としたが，構造化出力に従った\code{gpt-4.1-mini}/\code{gpt-5.4-mini}と完全に同じ運用条件とはみなさない．

表の数値は公開パッケージ内の最終ledgerを基準とし，\code{overall.csv}，\code{by\_stage.csv}，\code{by\_scenario\_group.csv}から作成した．

\end{document}
